% 第9章第4节：一般方阵特征值的扰动（Bauer--Fike界，特征基的条件，不可对角化的情形）。

\section{一般方阵特征值的扰动}
\label{sec:ls-nonnormal-spectral-sensitivity}

第~\ref{sec:ls-spectral-value-sensitivity}节按序比较特征值的界依靠对称结构：特征值是实数，可以排序，并能由二次型在特征子空间上的取值刻画。一般方阵的特征值可能是复数，没有自然的次序，二次型也不再与特征值对应。下面对可对角化的方阵研究另一个问题：$\bm A+\bm E$的某个特征值离$\bm A$的谱有多远。这里允许在复数域中选取特征基，向量使用复欧氏范数$\lVert\bm z\rVert_2=(\sum_i|z_i|^2)^{1/2}$，矩阵使用相应的诱导范数。

\subsection{Bauer--Fike界}

\begin{theorem}[Bauer--Fike]
\label{thm:ls-bauer-fike}
设$\bm A\in\mathbb C^{n\times n}$可对角化，$\bm A=\bm V\bm\Lambda\bm V^{-1}$，其中$\bm\Lambda=\operatorname{diag}(\lambda_1,\ldots,\lambda_n)$，$\bm V$可逆；记$\kappa_2(\bm V)=\lVert\bm V\rVert_2\lVert\bm V^{-1}\rVert_2$。则对任意$\bm E\in\mathbb C^{n\times n}$与$\bm A+\bm E$的任一特征值$\mu$，
\[
 \min_{1\leq i\leq n}|\mu-\lambda_i|\leq\kappa_2(\bm V)\lVert\bm E\rVert_2 .
\]
\end{theorem}
\begin{proof}
若$\mu$是$\bm A$的特征值，左端为零，结论成立。下设$\mu$不是$\bm A$的特征值，于是对角矩阵$\bm\Lambda-\mu\bm I$可逆。取$\bm A+\bm E$对应于$\mu$的特征向量$\bm w\neq\bm0$，由$(\bm A+\bm E)\bm w=\mu\bm w$得$(\bm A-\mu\bm I)\bm w=-\bm E\bm w$。代入$\bm A-\mu\bm I=\bm V(\bm\Lambda-\mu\bm I)\bm V^{-1}$，令$\bm y=\bm V^{-1}\bm w\neq\bm0$，得
\[
 \bm y=-(\bm\Lambda-\mu\bm I)^{-1}\bm V^{-1}\bm E\bm V\bm y .
\]
两边取范数，
\[
 \begin{aligned}
 \lVert\bm y\rVert_2
 &\leq\lVert(\bm\Lambda-\mu\bm I)^{-1}\rVert_2\,\lVert\bm V^{-1}\bm E\bm V\rVert_2\,\lVert\bm y\rVert_2\\
 &\leq\lVert(\bm\Lambda-\mu\bm I)^{-1}\rVert_2\,\kappa_2(\bm V)\lVert\bm E\rVert_2\,\lVert\bm y\rVert_2 .
 \end{aligned}
\]
对角矩阵$\bm D=\operatorname{diag}(d_i)$满足$\lVert\bm D\bm z\rVert_2^2=\sum_i|d_i|^2|z_i|^2\leq\max_i|d_i|^2\lVert\bm z\rVert_2^2$，且在相应的坐标向量处取等号，所以$\lVert(\bm\Lambda-\mu\bm I)^{-1}\rVert_2=1/\min_i|\lambda_i-\mu|$。代入并除以$\lVert\bm y\rVert_2>0$，得$1\leq\kappa_2(\bm V)\lVert\bm E\rVert_2/\min_i|\lambda_i-\mu|$，整理即得结论。
\end{proof}
界中出现的是所选特征基的条件数$\kappa_2(\bm V)$，而不是第~\ref{sec:ls-perturbation}节的$\kappa_2(\bm A)$，二者衡量不同的东西。$\bm A=\operatorname{diag}(1,s)$（$0<s<1$）可取$\bm V=\bm I$，$\kappa_2(\bm V)=1$，任何扰动都只能使特征值移动不超过$\lVert\bm E\rVert_2$，而$\kappa_2(\bm A)=1/s$可以很大。

\subsection{特征基的条件}
特征基的条件数与矩阵条件数可以相差很远。上面的对角矩阵是$\kappa_2(\bm A)$大而$\kappa_2(\bm V)$小的情形；反过来，求逆不敏感的矩阵也可能有敏感的特征值。取$\eta>0$与
\[
 \bm A=\begin{bmatrix}1&1\\0&1+\eta\end{bmatrix},\qquad
 \bm V=\begin{bmatrix}1&1\\0&\eta\end{bmatrix},
\]
由$\bm A(1,0)^\top=(1,0)^\top$与$\bm A(1,\eta)^\top=(1+\eta)(1,\eta)^\top$，$\bm V$的两列是特征值$1$与$1+\eta$的特征向量。矩阵条件数却不大：$\bm A^\top\bm A=\left[\begin{smallmatrix}1&1\\1&1+(1+\eta)^2\end{smallmatrix}\right]$的迹为$2+(1+\eta)^2$、行列式为$(1+\eta)^2$，$\eta\to0$时分别趋于$3$与$1$，它的两个特征值作为二次方程的根趋于$(3\pm\sqrt5)/2$，所以$\kappa_2(\bm A)$趋于$\sqrt{(3+\sqrt5)/(3-\sqrt5)}=(3+\sqrt5)/2\approx2.62$，求逆并不敏感。

但两个特征向量越来越接近平行：$\det(\bm V^\top\bm V)=(\det\bm V)^2$，所以$\bm V$的两个奇异值之积为$|\det\bm V|=\eta$，而最大奇异值不小于$\lVert\bm V\bm e_1\rVert_2=1$，故$\kappa_2(\bm V)\geq1/\eta$，Bauer--Fike界变得很宽。非正交的特征基会放大扰动在特征坐标中的表示，这是一般方阵特征值敏感的来源；对称矩阵可取正交的$\bm V$，$\kappa_2(\bm V)=1$，界与Weyl界同量级。

使用这个界时还要注意三点。第一，特征向量的缩放会改变$\bm V$的条件数，界对任意一组特征基成立，引用时须说明所用的基，通常取使$\kappa_2(\bm V)$较小的一组。第二，它只保证每个扰动后的特征值接近某个原特征值，不建立逐个特征值的固定配对，也不恢复对称情形的有序比较。第三，它要求$\bm A$可对角化；特征向量不足时，特征值的变化可以与扰动不成比例。

\subsection{不可对角化的情形}
特征向量不足时，下面的例子说明特征值的移动可以远大于扰动。
\begin{example}[特征值的移动远大于扰动]
\label{ex:ls-jordan-eigenvalue-sensitivity}
取不可对角化的
\[
 \bm A=\begin{bmatrix}0&1\\0&0\end{bmatrix},
\]
它的特征值全为$0$，特征空间只有一维。在左下角放一个大小为$t$的扰动$\bm E=\left[\begin{smallmatrix}0&0\\t&0\end{smallmatrix}\right]$，$t>0$，则$\lVert\bm E\rVert_2=t$，而
\[
 \bm A+\bm E=\begin{bmatrix}0&1\\t&0\end{bmatrix},\qquad
 \det\bigl(z\bm I-(\bm A+\bm E)\bigr)=z^2-t ,
\]
特征值为$\pm\sqrt t$，移动量为$\sqrt t$。扰动越小，这个比例反而越悬殊：$t=10^{-10}$时特征值移动$10^{-5}$，是扰动大小的$10^5$倍。由于$\sqrt t/t\to\infty$（$t\to0^+$），不存在常数$c$使移动量对所有小扰动都不超过$c\lVert\bm E\rVert_2$，Bauer--Fike那样与扰动成比例的界在这里不可能成立。若改取$t<0$，特征值为$\pm\mathrm i\sqrt{|t|}$，实特征值随之消失。
\end{example}
这个例子说明，缺少可对角化条件时，特征值的变化可以出现与扰动不同的量级；本书不为此引入新的标准形体系。第~\ref{sec:multivariable-eigenvalue-differential}节在相应光滑分支存在时研究一阶响应；本节的有限变化界则始终要连同对象的结构（对称、可对角化）和谱的分离条件一起使用。

特征值的位置是这样，特征方向又如何？下一节说明，一般方阵所选特征方向的偏转同样不能只由特征值间隔判断。
